\section{Introduction}
\label{sec:beer-introduction}

The preceding chapters have assembled, one by one, the ingredients of a verified translation: a set-theoretic universe (\Cref{ch:zflean}), a formal account of fragments of B and SMT-LIB with their respective semantics (\Cref{ch:b,ch:smt}), and a theory of representation changes between SMT types reconciling sets and types (\Cref{ch:loosening}).
This chapter puts them together.
It presents \beer{} (B Encoder), a tool written in Lean that translates Atelier~B proof obligations into SMT-LIB scripts, and---its actual subject---the mechanized proof that this translation is sound~\cite{trelat2025beer}.
The implementation is deliberately kept in the background: it is the topic of a single section (\Cref{sec:beer-implementation}), and everything before it is stated in ordinary mathematical language.

\subsection{What is proved}
\label{sec:beer-goal}

The correctness criterion for \beer{} is semantic equivalence: a B proof obligation and its SMT encoding must have the same truth value under corresponding valuations.
This is the end-to-end guarantee established in this chapter.

At the level of arbitrary terms, however, semantic equivalence cannot generally be stated as literal equality of denotations.
The translation changes representations: a B set becomes an SMT characteristic predicate, while a relation known to be functional may be represented by an option-valued function, as developed in \Cref{ch:loosening}.
Thus, a B term \(t\) of type \(\alpha\) and its SMT encoding \(t'\) of type \(\toSMT{\alpha}\) are denoted by elements of different carrier sets.
Their denotations must instead \emph{correspond} through the change of representation.

\begin{figure}[t]
  \begin{center}
    \begin{tikzpicture}
      \node (t) {$t$};
      \node (tsmt) [right=of t] {$\toSMT{t}$};
      \node (dent) [below=of t] {$T$};
      \node (dentsmt) [below=of tsmt] {$T'$};

      \draw[funarr] (t) -- (tsmt) node[midway, above] {\Large\textbeer};
      \draw[funarr] (t) -- (dent) node[midway, left] {$\interpB{\cdot}$};
      \draw[relarr, draw=black] (tsmt) -- (dentsmt) node[midway, right] {$\interpS{\cdot}$};
      \draw[relarr, draw=black] (dent) -- (dentsmt) node[midway, below] {$\semrel$};
    \end{tikzpicture}
  \end{center}
  \caption{A simplified view of the semantic square of \beer. Solid arrows are hypotheses, dashed arrows are conclusions.}
  \label{fig:beer-semantic-overview}
\end{figure}

The kernel of the correctness result is consequently an agreement statement.
We first state it here informally, just to motivate the rest of the chapter.
For every well-typed, well-defined B term \(t : \alpha\), the encoder produces an SMT term \(t' : \toSMT{\alpha}\) such that
\begin{equation*}
  \interpB{t} \semrel \interpS{t'},
  \label{eq:beer-invariant}
\end{equation*}
where \(\semrel\) is the agreement relation identifying an SMT denotation with the B denotation it represents.\footnote{The corresponding Unicode character \(\semrel\) is named \emph{corresponds to} (U+2258), which is its intended reading.}

Both denotations are sets in the ZFC universe of \Cref{ch:zflean} augmented with typing information (the semantic triples) as developed in \Cref{ch:b,ch:smt}: B terms and SMT terms denote in one and the same model, so agreement is a genuine set-theoretic relation---in fact, an identity up to a canonical change of representation---rather than transport between two unrelated semantics.
This idea ruled out the possibility of a semantic interpretation of SMT-LIB in bare higher-order logic (HOL), even though it would have been a very natural choice.
This design decision, made in \Cref{ch:b,ch:smt}, is what this chapter relies on.
\Cref{sec:beer-semantic-bridge} constructs \(\semrel\); \Cref{sec:beer-encoding-rules} gives the encoding rules; \Cref{sec:beer-soundness} proves the main correctness statement; and \Cref{sec:beer-proof-obligations} wraps it into the end-to-end guarantee for proof obligations.
The division between the executable encoder and this semantic argument is depicted in \Cref{fig:beer-architecture}, which also shows the overall architecture of \beer{}.

\begin{center}
\begin{minipage}{\textwidth}
  \centering
  \tikzsetnextfilename{beer-architecture}
  \begin{tikzpicture}[
    every node/.style={font=\small},
    node distance=14mm and 23mm,
    box/.style={
      draw=dark-blue,
      text=dark-blue,
      very thick,
      rounded corners=1mm,
      minimum width=19mm,
      minimum height=7mm,
      align=center,
    },
    concrete/.style={box, draw=dark-blue, text=dark-blue},
    phoas/.style={box, draw=orange, text=orange},
    semantic/.style={box, draw=pink, text=pink},
    note/.style={draw=none, inner sep=1.5pt},
    arr/.style={-{Stealth[length=2mm]}, very thick,},
    arr'/.style={-{Stealth[length=2mm]}, very thick, shorten >= 3pt, shorten <= 3pt},
    group/.style={
      draw=dark-blue!50,
      dashed,
      thick,
      rounded corners=1mm,
      inner sep=2mm,
    },
  ]
  \node[concrete] (b) {B};
  \node[concrete, right=of b] (smt) {SMT};

  \node[phoas, below=of b] (bphoas) {B PHOAS};
  \node[phoas, below=of smt] (smtphoas) {SMT PHOAS};
  \path (bphoas) -- coordinate[midway] (proofmid) (smtphoas);
  \node[semantic, below=12mm of proofmid] (zfc) {ZFC};

  \draw[arr', dark-blue] (b) -- node[midway, above]
    {\Large\textbeer} (smt);
  \draw[arr', orange] (b) -- node[midway, left, text=orange]
    {$\aterm{\cdot}{\Delta}$} (bphoas);
  \draw[arr', orange] (smt) -- node[midway, right, text=orange]
    {$\aterm{\cdot}{\Theta}$} (smtphoas);
  \draw[arr', pink] (bphoas) -- node[midway, below left, text=pink]
    {$\interpB{\cdot}$} (zfc);
  \draw[arr', pink] (smtphoas) -- node[midway, below right, text=pink]
    {$\interpS{\cdot}$} (zfc);

  \node[group, fit=(b)(smt)] (program) {};
  \node[note, text=dark-blue, font=\small,
    below=0pt of program.south] (programlabel) {program};

  \node[group, fit=(bphoas)(smtphoas)(zfc)] (proofs) {};
  \node[note, text=dark-blue, font=\small,
    above=0pt of proofs.north] (proofslabel) {proofs};

  \node[
    draw=pink,
    very thick,
    rounded corners=2mm,
    fit=(program)(programlabel)(proofs)(proofslabel),
    inner sep=2mm,
  ] (leanframe) {};
  \node[note, text=pink, font=\Large,
    above=0pt of leanframe.north] {\textlean};

  \node[box, left=8mm of leanframe] (pog) {POG file};
  \node[box, right=8mm of leanframe] (smtfile) {SMT file};

  \draw[very thick, draw=dark-blue, shorten <= 3pt] (pog.east) -- (leanframe.west);
  \draw[arr, draw=dark-blue, shorten >= 3pt] (leanframe.west) --
    node[midway, left, xshift=-2mm, text=dark-blue] {parse} (b.west);
  \draw[very thick, draw=dark-blue, shorten <= 3pt] (smt.east) --
    node[midway, right, xshift=2mm, text=dark-blue] {write} (leanframe.east);
  \draw[arr, draw=dark-blue, shorten >= 3pt] (leanframe.east) -- (smtfile.west);
\end{tikzpicture}

  \captionof{figure}{Architecture of \beer. The executable pipeline parses a POG file, encodes B as SMT, and writes an SMT-LIB file. The proof layer places both languages in PHOAS and compares their denotations in the common ZFC universe.}
  \label{fig:beer-architecture}
\end{minipage}
\end{center}

\subsection{A running example}
\label{sec:beer-running-example}

The following small example accompanies the whole chapter.
We consider two abstract sets \(X, Y\) and two constants
\[
  \var{f} \inB X \tfunB Y,
  \qquad
  \var{g} \inB X \mathbin{\bop{\leftrightarrow}} Y,
\]
a total function and a relation, and state the invariant \(\var{f} \cupB \var{g} \subseteqB X \prodB Y\).
The B method would generate one proof obligation for it, whose goal is the invariant itself under the two hypotheses above:
\begin{equation*}
  \bigl(\var{f} \inB X \tfunB Y\bigr) \andB
  \bigl(\var{g} \inB X \mathbin{\bop{\leftrightarrow}} Y\bigr)
  \impB
  \var{f} \cupB \var{g} \subseteqB X \prodB Y.
  \label{eq:beer-running-po}
\end{equation*}
The example is small but not innocent.
After elaboration of the derived forms of \Cref{ch:b} (\(X \bop{\leftrightarrow} Y\) is \(\powB(X \prodB Y)\), and \(X \tfunB Y\) a comprehension over \(X \pfunB Y\)), the two operands of \(\cupB\) live at the same B type but admit two different natural SMT representations: \(\var{f}\) is a partial \emph{function} value, \(\var{g}\) a \emph{relation} value.
Their union therefore exercises the multi-representation problem that motivated \Cref{ch:loosening}, and the inclusion goal---decoded as a bounded universal quantification---exercises the hardest case of the soundness proof.
The proof obligation \cref{eq:beer-running-po} is encoded in \Cref{sec:beer-encoding-rules}, its soundness instance is identified in \Cref{sec:beer-soundness}, and the complete emitted script is shown in \Cref{sec:beer-proof-obligations}.
